% ============================================================
%  يُترجَم بـ XeLaTeX  (في Overleaf: Menu → Compiler → XeLaTeX)
%  الخط: Amiri (متوفر في Overleaf، أو حمّله من github.com/aliftype/amiri)
% ============================================================
\documentclass[12pt,a4paper]{article}
\usepackage[a4paper,margin=1.7cm,top=2.3cm,bottom=2.9cm,headheight=15pt,footskip=1.5cm]{geometry}
\usepackage{amsmath,amssymb}
\usepackage{icomma}
\usepackage[english,bidi=default,layout=lists]{babel}
\babelprovide[import,main]{arabic}
\babelfont{rm}[Scale=1.08]{Amiri}
\babelfont{sf}[Scale=1.08]{Amiri}
\usepackage[most]{tcolorbox}
\usepackage{tikz}
\usetikzlibrary{arrows.meta,calc,shadings}
\usepackage{array}
\usepackage{colortbl}
\usepackage{fancyhdr}
\usepackage{titlesec}
\usepackage{enumitem}
\usepackage{collcell}
\usepackage{needspace}
\usepackage{etoolbox}
% أعمدة جداول تُكتب خلاياها بالعربية في الاتجاه الصحيح
\newcolumntype{C}{>{\collectcell\ar}c<{\endcollectcell}}
% تجنب عنوان يتيم في أسفل الصفحة
\pretocmd{\section}{\Needspace{15\baselineskip}}{}{}
\pretocmd{\subsection}{\Needspace{11\baselineskip}}{}{}

% ---------- الألوان ----------
\definecolor{cDef}{HTML}{1F4E9C}
\definecolor{cThm}{HTML}{B3261E}
\definecolor{cRem}{HTML}{2E7D32}
\definecolor{cEx}{HTML}{C76B00}
\definecolor{cMeth}{HTML}{6A1B9A}
\definecolor{cSol}{HTML}{00695C}
\definecolor{cSec}{HTML}{1F3A68}

% ---------- العناوين ----------
\titleformat{\section}{\Large\bfseries\color{cSec}}{\thesection.}{0.6em}{}[{\color{cSec!60}\titlerule[1pt]}]
\titleformat{\subsection}{\large\bfseries\color{cThm!85!black}}{\thesubsection.}{0.5em}{}
\titlespacing*{\section}{0pt}{14pt}{8pt}
\titlespacing*{\subsection}{0pt}{10pt}{5pt}
\setlength{\parindent}{0pt}
\setlength{\parskip}{4pt}
\setlist{itemsep=2pt,topsep=3pt}

% ---------- الصناديق (تُرسم في سياق LTR ويُكتب محتواها بالعربية) ----------
\newtcolorbox{arbox}[3][]{enhanced,
  colback=#2!5,colframe=#2!80!black,arc=2mm,boxrule=0.8pt,
  left=6pt,right=6pt,top=8pt,bottom=5pt,
  fonttitle=\bfseries,coltitle=white,
  title={\foreignlanguage{arabic}{#3}},
  attach boxed title to top right={xshift=-5mm,yshift=-3mm},
  boxed title style={colback=#2!80!black,arc=1.5mm,boxrule=0pt},
  before upper={\selectlanguage{arabic}},
  before skip=8pt,after skip=8pt,breakable,#1}
\newenvironment{BOX}[2]{\begin{otherlanguage}{english}\begin{arbox}{#1}{#2}}{\end{arbox}\end{otherlanguage}}
\newenvironment{tarif}[1][تعريف]{\begin{BOX}{cDef}{#1}}{\end{BOX}}
\newenvironment{khassia}[1][خاصية]{\begin{BOX}{cThm}{#1}}{\end{BOX}}
\newenvironment{molahada}[1][ملاحظة]{\begin{BOX}{cRem}{#1}}{\end{BOX}}
\newenvironment{mithal}[1][مثال]{\begin{BOX}{cEx}{#1}}{\end{BOX}}
\newenvironment{tariqa}[1][طريقة]{\begin{BOX}{cMeth}{#1}}{\end{BOX}}

% ---------- التمارين والحلول ----------
\newcounter{exo}
\newenvironment{tamrin}[1]{\stepcounter{exo}\begin{BOX}{cSec}{التمرين \theexo: #1}}{\end{BOX}}
\newenvironment{hal}{\Needspace{12\baselineskip}\begin{otherlanguage}{english}\begin{arbox}[breakable]{cSol}{حل التمرين \theexo}}{\end{arbox}\end{otherlanguage}}

% رسومات TikZ تُوضع في سياق LTR
\newenvironment{LTR}{\begin{otherlanguage}{english}}{\end{otherlanguage}}
\newcommand{\ar}[1]{\foreignlanguage{arabic}{#1}}

% ---------- رموز ----------
% نقاط الحذف داخل الرياضيات (\dots يسبب خللاً في الاتجاه مع babel)
\renewcommand{\dots}{\mathinner{\ldotp\ldotp\ldotp}}
\renewcommand{\ldots}{\mathinner{\ldotp\ldotp\ldotp}}
\renewcommand{\cdots}{\mathinner{\cdotp\cdotp\cdotp}}
\newcommand{\N}{\mathbb{N}}
\newcommand{\Z}{\mathbb{Z}}
\newcommand{\D}{\mathbb{D}}
\newcommand{\Q}{\mathbb{Q}}
\newcommand{\R}{\mathbb{R}}
\newcommand{\pgcd}{\operatorname{PGCD}}
\newcommand{\ppcm}{\operatorname{PPCM}}

% ---------- تمثيل مجال على مستقيم عددي ----------
% \NLseg{a}{b}{[}{]}  مجال محدود
\newcommand{\NLseg}[4]{\begin{LTR}\begin{tikzpicture}[baseline=-0.6ex]
 \draw[-{Stealth[length=2mm]},gray!70](0,0)--(3.4,0);
 \draw[line width=2pt,cThm!80!black](0.7,0)--(2.7,0);
 \node[cThm!80!black,font=\LARGE] at (0.7,0){$#3$};
 \node[cThm!80!black,font=\LARGE] at (2.7,0){$#4$};
 \node[below=3pt,font=\footnotesize] at (0.7,0){$#1$};
 \node[below=3pt,font=\footnotesize] at (2.7,0){$#2$};
\end{tikzpicture}\end{LTR}}
% \NLright{a}{[}  مجال نحو +∞
\newcommand{\NLright}[2]{\begin{LTR}\begin{tikzpicture}[baseline=-0.6ex]
 \draw[-{Stealth[length=2mm]},gray!70](0,0)--(3.4,0);
 \draw[line width=2pt,cThm!80!black](1.3,0)--(3.3,0);
 \node[cThm!80!black,font=\LARGE] at (1.3,0){$#2$};
 \node[below=3pt,font=\footnotesize] at (1.3,0){$#1$};
\end{tikzpicture}\end{LTR}}
% \NLleft{b}{]}  مجال نحو −∞
\newcommand{\NLleft}[2]{\begin{LTR}\begin{tikzpicture}[baseline=-0.6ex]
 \draw[-{Stealth[length=2mm]},gray!70](0,0)--(3.4,0);
 \draw[line width=2pt,cThm!80!black](0.05,0)--(2.1,0);
 \node[cThm!80!black,font=\LARGE] at (2.1,0){$#2$};
 \node[below=3pt,font=\footnotesize] at (2.1,0){$#1$};
\end{tikzpicture}\end{LTR}}

% ---------- أيقونات مواقع التواصل (مرسومة بـ TikZ، لا تحتاج أي حزمة) ----------
\definecolor{cFB}{HTML}{1877F2}
\definecolor{cTG}{HTML}{229ED9}
\definecolor{cIG1}{HTML}{FEDA75}
\definecolor{cIG2}{HTML}{D62976}
\definecolor{cIG3}{HTML}{4F5BD5}
\newcommand{\fbicon}{\begin{LTR}\begin{tikzpicture}[baseline=-0.75ex,x=1em,y=1em]
 \fill[cFB] (0,0) circle (0.55);
 \node[white,inner sep=0,font=\usefont{TU}{lmss}{bx}{n}\small] at (0.04,-0.02) {f};
\end{tikzpicture}\end{LTR}}
\newcommand{\igicon}{\begin{LTR}\begin{tikzpicture}[baseline=-0.75ex,x=1em,y=1em]
 \shade[lower left=cIG1,lower right=cIG2,upper left=cIG2,upper right=cIG3,rounded corners=0.28em] (-0.55,-0.55) rectangle (0.55,0.55);
 \draw[white,line width=0.09em,rounded corners=0.18em] (-0.36,-0.36) rectangle (0.36,0.36);
 \draw[white,line width=0.09em] (0,0) circle (0.17);
 \fill[white] (0.22,0.22) circle (0.05);
\end{tikzpicture}\end{LTR}}
\newcommand{\tgicon}{\begin{LTR}\begin{tikzpicture}[baseline=-0.75ex,x=1em,y=1em]
 \fill[cTG] (0,0) circle (0.55);
 \fill[white] (-0.33,0.0) -- (0.32,0.27) -- (0.2,-0.3) -- (0.0,-0.15) -- (-0.08,-0.28) -- (-0.08,-0.08) -- cycle;
 \fill[cTG!70!black] (-0.08,-0.08) -- (0.25,0.2) -- (0.0,-0.15) -- (-0.08,-0.28) -- cycle;
\end{tikzpicture}\end{LTR}}
% سطر مواقع التواصل أسفل كل صفحة
% (يُكتب السطر بترتيب بصري من اليسار إلى اليمين: الأيقونة يمين كل اسم)
\newcommand{\PageBox}{\begin{LTR}\tikz[baseline=(p.base)]\node[fill=black,text=white,inner sep=2.5pt,minimum width=1.6em,minimum height=1.3em,font=\usefont{TU}{lmss}{bx}{n}\small](p){\ifnum\value{page}<10 0\fi\arabic{page}};\end{LTR}}
\newcommand{\SocialRow}{\begin{otherlanguage}{english}%
 \ar{\textbf{الأستاذ سوايسية هشام للرياضيات}}\ \tgicon\hfill
 \textbf{prof\_hicham\_math}\ \igicon\hfill
 \ar{\textbf{الأستاذ سوايسية هشام للرياضيات}}\ \fbicon\hspace{2.2em}\PageBox\end{otherlanguage}}
% ---------- رأس وتذييل الصفحة ----------
\pagestyle{fancy}
\fancyhf{}
\renewcommand{\headrulewidth}{0.6pt}
\renewcommand{\footrulewidth}{0.4pt}
% \SetHeader{يمين الرأس}{يسار الرأس}{يمين التذييل}{يسار التذييل}
\newcommand{\SetHeader}[2]{%
 \fancyhead[C]{\parbox{\textwidth}{\ar{#1\hfill #2}}}%
 \fancyfoot[C]{\parbox{\textwidth}{\small\SocialRow}}}

\usepackage{tkz-tab}
\usepackage{multicol}
% جدول تغيرات / إشارة (يُرسم في سياق LTR)
\newenvironment{TAB}{\par\begin{center}\begin{LTR}\begin{tikzpicture}[scale=0.95,every node/.style={font=\small}]}{\end{tikzpicture}\end{LTR}\end{center}}
% معلم بشبكة: \begin{GR}{xmin}{xmax}{ymin}{ymax}{unit cm}
\newenvironment{GR}[5]{\par\begin{center}\begin{LTR}\begin{tikzpicture}[x=#5cm,y=#5cm,font=\footnotesize]
 \def\xmn{#1}\def\xmx{#2}\def\ymn{#3}\def\ymx{#4}
 \draw[gray!25,very thin] (#1,#3) grid (#2,#4);
 \draw[-{Stealth[length=2mm]}] (#1,0)--(#2+0.3,0) node[below left]{$x$};
 \draw[-{Stealth[length=2mm]}] (0,#3)--(0,#4+0.3) node[below left]{$y$};
 \node[below left] at (0,0){$O$};
 \draw[very thick,-{Stealth[length=1.6mm]}] (0,0)--(1,0) node[below,pos=.55]{$\vec i$};
 \draw[very thick,-{Stealth[length=1.6mm]}] (0,0)--(0,1) node[left,pos=.55]{$\vec j$};
 \begin{scope}\clip (#1,#3) rectangle (#2,#4);}
{\end{scope}\end{tikzpicture}\end{LTR}\end{center}}
\newcommand{\Ticks}[2]{\foreach \t in {#1} \node[below,font=\scriptsize] at (\t,0){$\t$}; \foreach \t in {#2} \node[left,font=\scriptsize] at (0,\t){$\t$};}
\newcommand{\bx}[1]{\ensuremath{\colorbox{cDef!12}{$\displaystyle #1$}}}
\newcommand{\hl}[1]{\ensuremath{\colorbox{cEx!20}{$\displaystyle #1$}}}
\newcommand{\Lim}[1]{\displaystyle\lim_{x\to #1}}
\newcommand{\Cf}{(C_f)}
\newenvironment{qst}{\begin{enumerate}[label=\protect\circled{\arabic*},leftmargin=*]}{\end{enumerate}}
\newcommand{\circled}[1]{\foreignlanguage{english}{\tikz[baseline=(c.base)]\node[circle,fill=cDef,text=white,inner sep=1.2pt,font=\scriptsize\bfseries](c){#1};}}
% عمودان: \TwoCol{عرض اليسار}{محتوى اليسار}{عرض اليمين}{محتوى اليمين}  (آمن مع TikZ)
\newcommand{\TwoCol}[4]{\par\begin{LTR}\noindent\begin{minipage}[t]{#1\textwidth}\selectlanguage{arabic}#2\end{minipage}\hfill\begin{minipage}[t]{#3\textwidth}\selectlanguage{arabic}#4\end{minipage}\end{LTR}\par}
\newcommand{\un}{(u_n)}\newcommand{\vn}{(v_n)}
% برهان بالتراجع: \Rec{الخاصية P(n)}{التحقق من P(0)}{البرهان على الوراثة}
\newcommand{\Rec}[3]{\textbf{برهان بالتراجع:} نسمي $P(n)$ الخاصية: «#1».
\begin{itemize}
\item \textbf{التحقق:} #2
\item \textbf{الوراثة:} نفرض أن $P(n)$ صحيحة ونبرهن صحة $P(n+1)$: #3
\item \textbf{الخلاصة:} حسب مبدأ الاستدلال بالتراجع، $P(n)$ صحيحة من أجل كل عدد طبيعي $n$.
\end{itemize}}
\newcommand{\src}[1]{{\footnotesize\textbf{#1}}\par}
\newcommand{\ok}{$\textcolor{cRem}{\checkmark}$}
\newcommand{\limn}{\displaystyle\lim_{n\to+\infty}}
% أسئلة فرعية (أ) (ب) (جـ) كل سؤال في سطر
\newenvironment{sq}{\begin{itemize}[leftmargin=2.4em,itemsep=3pt]}{\end{itemize}}
\newcommand{\qa}{\item[(أ)]}\newcommand{\qb}{\item[(ب)]}\newcommand{\qc}{\item[(جـ)]}\newcommand{\qd}{\item[(د)]}\newcommand{\qe}{\item[(هـ)]}
\newcommand{\pa}{\par\smallskip\textbf{(أ)} }\newcommand{\pb}{\par\smallskip\textbf{(ب)} }\newcommand{\pc}{\par\smallskip\textbf{(جـ)} }\newcommand{\pd}{\par\smallskip\textbf{(د)} }\newcommand{\pe}{\par\smallskip\textbf{(هـ)} }
\newcommand{\ccl}[1]{\par\textbf{الخلاصة:} #1}

\pretocmd{\section}{\Needspace{24\baselineskip}}{}{}
\SetHeader{\small\bfseries\color{cSec}تمارين: الجداء السلمي في المستوي}{\small\color{cSec}السنة الثانية ثانوي -- الشعب العلمية}
\begin{document}
\newcolumntype{M}[1]{>{\centering\arraybackslash}p{#1}}
\newcommand{\RL}[4]{\TwoCol{#3}{#4}{#1}{#2}}
\newcommand{\QQ}[1]{\par\smallskip\textbf{#1.} }
\newcommand{\V}[1]{\overrightarrow{#1}}
\newcommand{\nV}[1]{\big\|#1\big\|}
\newcommand{\ABC}[3]{\\[2pt]\begin{tabular}{p{0.28\textwidth}p{0.28\textwidth}p{0.28\textwidth}}(جـ) \ #3 & (ب) \ #2 & (أ) \ #1\end{tabular}}
\begin{LTR}
\begin{tcolorbox}[enhanced,colback=cSec,colframe=cSec,arc=3mm,halign=center,boxrule=0pt,top=12pt,bottom=12pt,before upper={\selectlanguage{arabic}}]
{\color{white}\fontsize{26}{32}\selectfont\bfseries سلسلة تمارين: الجداء السلمي في المستوي}\\[8pt]
{\color{white}\large تمارين متنوعة ومسائل شاملة -- بدون حلول (للتدرّب)}\\[5pt]
{\color{white}\large السنة الثانية ثانوي -- الشعب العلمية}
\end{tcolorbox}
\end{LTR}

% =====================================================================
\section{حساب الجداء السلمي والعلاقات المترية}
% ------------------------------------------------- 1
\begin{tamrin}{مثلث متقايس الأضلاع}
$ABC$ مثلث متقايس الأضلاع طول ضلعه $4$ ، و $I$ منتصف القطعة $[BC]$.
\begin{enumerate}
\item احسب الجداءات السلمية:
\begin{sq}
\qa $\V{AB}\cdot\V{AC}$
\qb $\V{BA}\cdot\V{BC}$
\qc $\V{AI}\cdot\V{BC}$
\qd $\V{AB}\cdot\V{AI}$
\end{sq}
\item احسب $\nV{\V{AB}+\V{AC}}$ ، ثم استنتج الطول $AI$.
\end{enumerate}
\end{tamrin}

% ------------------------------------------------- 2
\begin{tamrin}{مبرهنة الكاشي وقانون الجيوب}
$ABC$ مثلث حيث: $AB=5$ ، $AC=8$ و $\widehat{BAC}=\dfrac{\pi}{3}$. \ $I$ منتصف $[BC]$.
\begin{enumerate}
\item احسب $\V{AB}\cdot\V{AC}$ ، ثم الطول $BC$.
\item احسب مساحة المثلث $ABC$.
\item احسب $\cos\widehat{ABC}$ ثم $\sin\widehat{ABC}$.
\item تحقق من قانون الجيوب ، ثم استنتج نصف قطر الدائرة المحيطة بالمثلث $ABC$.
\item احسب طول المتوسط $AI$.
\end{enumerate}
\end{tamrin}

% ------------------------------------------------- 3
\begin{tamrin}{مجموعات نقط على قطعة}
لتكن $A$ و $B$ نقطتين متمايزتين من المستوي حيث $AB=5$ ، و $I$ منتصف $[AB]$.
\begin{enumerate}
\item نعتبر النقطة $C$ حيث $3\V{CA}+2\V{CB}=\vec0$.
\begin{sq}
\qa اكتب $\V{AC}$ بدلالة $\V{AB}$.
\qb لتكن $(E)$ مجموعة النقط $M$ من المستوي التي تحقق $\V{AM}\cdot\V{AB}=10$. \ بيّن أن $C\in(E)$ ، ثم عيّن وأنشئ $(E)$.
\end{sq}
\item $(F)$ مجموعة النقط $M$ التي تحقق $MA^2+MB^2=7$. \ عيّن $(F)$.
\item $(\Gamma)$ مجموعة النقط $M$ التي تحقق $\big(\V{MA}+\V{MB}\big)\cdot\V{MA}=0$. \ عيّن ثم أنشئ $(\Gamma)$.
\end{enumerate}
\end{tamrin}

% ------------------------------------------------- 4
\begin{tamrin}{مجموعة النقط $MB=2MA$}
لتكن $A$ و $B$ نقطتين متمايزتين حيث $AB=6$ ، و $(E)$ مجموعة النقط $M$ من المستوي التي تحقق: $MB=2MA$.
\begin{enumerate}
\item برهن أنه من أجل كل نقطة $M$: \ $MB=2MA\iff\big(\V{MB}+2\V{MA}\big)\cdot\big(\V{MB}-2\V{MA}\big)=0$.
\item لتكن $I$ و $J$ نقطتين حيث: $2\V{IA}+\V{IB}=\vec0$ و $-2\V{JA}+\V{JB}=\vec0$.
\begin{sq}
\qa أنشئ النقطتين $I$ و $J$.
\qb أثبت أنه من أجل كل نقطة $M$: \ $\big(\V{MB}+2\V{MA}\big)\cdot\big(\V{MB}-2\V{MA}\big)=-3\,\V{MI}\cdot\V{MJ}$.
\end{sq}
\item عيّن ثم أنشئ المجموعة $(E)$.
\end{enumerate}
\end{tamrin}

% ------------------------------------------------- 5
\begin{tamrin}{مجموعة النقط $MA^2-MB^2=k$}
في معلم متعامد ومتجانس نعتبر النقطتين $A(1;0)$ و $B(4;0)$ ، فيكون $AB=3$.
\begin{enumerate}
\item بيّن أنه من أجل كل نقطة $M$: \ $MA^2-MB^2=2\,\V{AM}\cdot\V{AB}-9$.
\item عيّن مجموعة النقط $M$ التي تحقق: $MA^2-MB^2=-3$ ، ثم اكتب معادلة لها.
\end{enumerate}
\end{tamrin}

% =====================================================================
\section{المستقيمات والدوائر في معلم متعامد ومتجانس}
% ------------------------------------------------- 6
\begin{tamrin}{دائرة ومماس}
المستوي منسوب إلى معلم متعامد ومتجانس $(O;\vec i,\vec j)$. نعتبر النقط: $A(1;-1)$ ، $B(1;3)$ ، $C(5;2)$ و $D(0;-2)$.
\begin{enumerate}
\item احسب مركبات الشعاعين $\V{AB}$ و $\V{AC}$ ، ثم احسب $\V{AB}\cdot\V{AC}$.
\item احسب $AB$ و $AC$ ثم استنتج $\cos\big(\V{AB};\V{AC}\big)$.
\item لتكن $(c)$ مجموعة النقط $M(x;y)$ التي تحقق: $x^2+y^2-2x-6y-3=0$. \ بيّن أن $(c)$ دائرة يطلب تعيين مركزها $\Omega$ ونصف قطرها $r$.
\item جد معادلة للمستقيم $(\Delta)$ الذي يشمل $D$ و $\vec n(3;2)$ شعاع ناظمي له.
\item
\begin{sq}
\qa بيّن أن المسافة بين $\Omega$ والمستقيم $(\Delta)$ هي $\sqrt{13}$.
\qb هل المستقيم $(\Delta)$ مماس للدائرة $(c)$؟ برّر.
\end{sq}
\end{enumerate}
\end{tamrin}

% ------------------------------------------------- 7
\begin{tamrin}{مثلث قائم ودائرة}
المستوي منسوب إلى معلم متعامد ومتجانس $(O;\vec i,\vec j)$. لتكن النقط $A(-1;1)$ ، $B(-2;1)$ و $C(-1;-1)$.
\begin{enumerate}
\item احسب الجداء السلمي $\V{AB}\cdot\V{AC}$ ، ثم استنتج طبيعة المثلث $ABC$.
\item أوجد معادلة للمستقيم $(\Delta)$ الذي يشمل $A$ و $\V{BC}$ شعاع ناظمي له.
\item أوجد معادلة الدائرة $(c)$ التي مركزها $w(-1;3)$ ونصف قطرها $r=\sqrt5$.
\item تحقق أن $B$ تنتمي إلى الدائرة $(c)$.
\item احسب المسافة بين مركز الدائرة $(c)$ والمستقيم $(\Delta)$ ، ثم استنتج الوضع النسبي لـ $(c)$ و $(\Delta)$.
\end{enumerate}
\end{tamrin}

% ------------------------------------------------- 8
\begin{tamrin}{الدائرة المحيطة بمثلث قائم وصورتها بتحاكٍ}
في معلم متعامد ومتجانس نعتبر النقط $A(2;0)$ ، $B(4;2)$ ، $C(-1;3)$ والمستقيم $(d):x+5y-21=0$.
\begin{enumerate}
\item احسب الجداء السلمي $\V{AB}\cdot\V{AC}$ ثم استنتج طبيعة المثلث $ABC$.
\item عيّن إحداثيي النقطة $D$ مركز الدائرة $(C)$ المحيطة بالمثلث $ABC$.
\item اكتب معادلة ديكارتية للدائرة $(C)$.
\item اكتب معادلة ديكارتية للمستقيم $(\Delta)$ المماس للدائرة $(C)$ في النقطة $B$ ، ثم تحقق أن $(\Delta)\perp(d)$.
\item احسب المسافة بين $D$ والمستقيم $(d)$ ، ثم ادرس وضعية $(d)$ بالنسبة إلى $(C)$.
\item عيّن صورة الدائرة $(C)$ بالتحاكي $h$ الذي مركزه $B$ ونسبته $2$.
\end{enumerate}
\end{tamrin}

% ------------------------------------------------- 9
\begin{tamrin}{تقاطع مستقيم ودائرة ومساحة مثلث}
في معلم متعامد ومتجانس نعتبر النقطتين $A(1;2)$ ، $B(0;5)$ والدائرة $(C):x^2+y^2-2x-3=0$.
\begin{enumerate}
\item اكتب معادلة ديكارتية للمستقيم $(D)$ الذي يشمل $I(-1;1)$ و $\vec u(2;1)$ شعاع توجيه له.
\item حدّد المركز $\Omega$ ونصف القطر $r$ للدائرة $(C)$ ، ثم تحقق أن $A\in(C)$.
\item اكتب معادلة ديكارتية للمستقيم $(\Delta)$ المار من $B$ و $\vec n(3;4)$ شعاع ناظمي له.
\item بيّن أن تقاطع $(\Delta)$ مع $(C)$ مجموعة خالية.
\item تحقق أن $(C)$ و $(D)$ يتقاطعان في نقطتين ، الأولى هي $A$ والثانية نرمز لها بـ $F$.
\item احسب مساحة المثلث $A\Omega F$ إذا علمت أن $\sin\widehat{A\Omega F}=0{,}8$.
\item احسب المسافة بين $\Omega$ والمستقيم $(D)$ بطريقتين.
\end{enumerate}
\end{tamrin}

% ------------------------------------------------- 10
\begin{tamrin}{دائرة ومماسان متوازيان}
المستوي منسوب إلى معلم متعامد ومتجانس $(O;\vec i,\vec j)$. نعتبر النقط $B\big(0;\sqrt3\big)$ ، $C\big(0;-\sqrt3\big)$ و $M(x;y)$.
\begin{enumerate}
\item بيّن أن مجموعة النقط $M$ من المستوي التي تحقق $\V{MB}\cdot\V{MC}=0$ هي الدائرة $(C)$ ، يطلب تحديد مركزها ونصف قطرها.
\item اكتب معادلة المستقيمين $(T_1)$ و $(T_2)$ مماسي الدائرة $(C)$ في النقطتين $C$ و $B$ على الترتيب.
\item بيّن أن $(T_1)$ يوازي $(T_2)$.
\item ارسم الدائرة $(C)$ والمستقيمين $(T_1)$ و $(T_2)$.
\end{enumerate}
\end{tamrin}

% ------------------------------------------------- 11
\begin{tamrin}{مستقيمان متعامدان ومجموعة نقط}
في معلم متعامد ومتجانس $(O;\vec i,\vec j)$ ، ليكن $(D)$ المستقيم الذي معادلته: $3x+3y-5=0$.
\begin{enumerate}
\item اكتب معادلة للمستقيم $(\Delta)$ الذي يشمل $A(-1;2)$ ويعامد $(D)$.
\item عيّن إحداثيي النقطة $B$ نقطة تقاطع المستقيمين $(D)$ و $(\Delta)$ ، ثم احسب $AB$.
\item عيّن مجموعة النقط $M(x;y)$ من المستوي التي تحقق: $\V{MA}\cdot\V{MB}=AB$.
\end{enumerate}
\end{tamrin}

% ------------------------------------------------- 12
\begin{tamrin}{دائرة ، مماس ونظيرة}
المستوي منسوب إلى معلم متعامد ومتجانس $(O;\vec i,\vec j)$. لتكن $(C)$ مجموعة النقط $M(x;y)$ المعرّفة بـ: $\dfrac12\big(x^2+y^2\big)-3x+y+\dfrac52=0$.
\begin{enumerate}
\item
\begin{sq}
\qa بيّن أن $(C)$ دائرة يطلب تحديد إحداثيي مركزها $\Omega$ ونصف قطرها.
\qb أثبت أن النقطة $A(1;-2)$ من $(C)$.
\end{sq}
\item اكتب معادلة المستقيم $(T)$ مماس الدائرة $(C)$ عند النقطة $A$.
\item لتكن $B$ نظيرة $\Omega$ بالنسبة إلى $A$.
\begin{sq}
\qa عيّن إحداثيي النقطة $B$.
\qb اكتب معادلة الدائرة $(C')$ التي مركزها $B$ والمماسة لـ $(C)$ عند النقطة $A$.
\end{sq}
\item ادرس الوضع النسبي بين المستقيم $(D):6x+2y-5=0$ والدائرة $(C)$.
\end{enumerate}
\end{tamrin}

% ------------------------------------------------- 13
\begin{tamrin}{مسألة شاملة: مماس ، جداء سلمي ومجموعات نقط}
المستوي منسوب إلى معلم متعامد ومتجانس $(O;\vec i,\vec j)$ حيث $\nV{\vec i}=\nV{\vec j}=1\,cm$. \ لتكن النقط $A(2;1)$ ، $B(5;-1)$ و $C(8;3)$ ، و $H$ منتصف القطعة $[AC]$.
\begin{enumerate}
\item علّم النقط $A$ ، $B$ ، $C$ و $H$ ، ثم اكتب معادلة ديكارتية للمستقيم $(\Delta)$ الذي يشمل $A$ و $\vec v(3;-2)$ ناظمي له.
\item لتكن $(\gamma)$ مجموعة النقط $M(x;y)$ حيث: $x^2+y^2-10x+2y+13=0$.
\begin{sq}
\qa أثبت أن $(\gamma)$ دائرة مركزها $B$ ونصف قطرها $R$ يطلب تعيينه.
\qb أنشئ المستقيم $(\Delta)$ والدائرة $(\gamma)$.
\qc تحقق حسابياً أن $A\in(\gamma)$ ، ثم حدّد بدقة المسافة $d(B;\Delta)$.
\qd استنتج الوضعية النسبية لـ $(\Delta)$ و $(\gamma)$.
\end{sq}
\item احسب الجداء السلمي $\V{BA}\cdot\V{BC}$ بطريقتين ، واستنتج قيمة مقرّبة بالدرجات لقيس الزاوية $\widehat{ABC}$.
\item احسب الطول $BH$ بطريقتين مختلفتين.
\item حدّد طبيعة وعناصر $(E)$ مجموعة النقط $N$ من المستوي التي تحقق: $NA^2+NC^2=21$.
\item حدّد طبيعة وعناصر $(F)$ مجموعة النقط $M$ التي تحقق: $\V{MA}\cdot\V{MC}=0$ ، واكتب معادلة ديكارتية لها.
\item ليكن $h$ التحاكي الذي مركزه $O(0;0)$ ونسبته $-\dfrac12$ ، و $(\gamma')$ صورة الدائرة $(\gamma)$ بالتحاكي $h$.
\begin{sq}
\qa اكتب معادلة ديكارتية للدائرة $(\gamma')$.
\qb استنتج محيط ومساحة الدائرة $(\gamma')$.
\end{sq}
\end{enumerate}
\end{tamrin}

% =====================================================================
\section{دساتير الجمع والمعادلات المثلثية}
% ------------------------------------------------- 14
\begin{tamrin}{حساب قيم مضبوطة}
\begin{enumerate}
\item إذا علمت أن $\sin\Big(\dfrac{7\pi}{12}\Big)=\dfrac{\sqrt2+\sqrt6}{4}$ ، احسب $\sin\Big(\dfrac{5\pi}{12}\Big)$ بطريقتين مختلفتين.
\item احسب $\cos\Big(\dfrac{5\pi}{12}\Big)$ ، ثم استنتج أن $\tan\Big(\dfrac{5\pi}{12}\Big)=2+\sqrt3$.
\item باستعمال دستور التضعيف ، احسب $\cos\Big(\dfrac{\pi}{8}\Big)$ و $\sin\Big(\dfrac{\pi}{8}\Big)$.
\end{enumerate}
\end{tamrin}

% ------------------------------------------------- 15
\begin{tamrin}{متطابقات مثلثية}
برهن أنه من أجل كل عدد حقيقي $x$:
\begin{sq}
\qa $\cos^4x-\sin^4x=\cos2x$
\qb $\big(\sin x+\cos x\big)^2=1+\sin2x$
\qc $\sin3x=3\sin x-4\sin^3x$
\qd $\cos3x=4\cos^3x-3\cos x$
\end{sq}
\end{tamrin}

% ------------------------------------------------- 16
\begin{tamrin}{معادلات من الشكل $a\cos x+b\sin x=c$}
\begin{enumerate}
\item حل في المجال $]-\pi;\pi]$ المعادلة: $\cos x+\sin x=1$.
\item حل في $\R$ المعادلة: $\sqrt3\cos x-\sin x=\sqrt2$.
\item حل في $\R$ المعادلة: $\cos4x=\dfrac{1}{\sqrt2}$.
\end{enumerate}
\end{tamrin}

% =====================================================================
\section{اختيار من متعدد ، صحيح أم خطأ}
% ------------------------------------------------- 17
\begin{tamrin}{اختيار من متعدد}
اختر الإجابة الصحيحة الوحيدة مع التعليل (المعلم متعامد ومتجانس):
\begin{enumerate}
\item إذا كان $\vec u(2;-1)$ و $\vec v(3;4)$ فإن $\vec u\cdot\vec v$ يساوي:
\ABC{$2$}{$10$}{$-2$}
\item $\nV{\vec u}=3$ ، $\nV{\vec v}=2$ و $(\vec u;\vec v)=\dfrac{2\pi}{3}$ ، إذن $\vec u\cdot\vec v$ يساوي:
\ABC{$3$}{$-3$}{$-3\sqrt3$}
\item المسافة بين المبدأ $O$ والمستقيم $3x-4y+10=0$ هي:
\ABC{$10$}{$2$}{$\dfrac{10}{7}$}
\item الدائرة $x^2+y^2+4x-6y+4=0$ مركزها ونصف قطرها:
\ABC{$(2;-3)$ ، $3$}{$(-2;3)$ ، $9$}{$(-2;3)$ ، $3$}
\item $\cos\dfrac{\pi}{12}$ يساوي:
\ABC{$\dfrac{\sqrt6+\sqrt2}{4}$}{$\dfrac{\sqrt6-\sqrt2}{4}$}{$\dfrac{\sqrt3}{4}$}
\item $A$ و $B$ نقطتان متمايزتان. مجموعة النقط $M$ حيث $\V{MA}\cdot\V{MB}=0$ هي:
\ABC{محور $[AB]$}{المستقيم $(AB)$}{الدائرة ذات القطر $[AB]$}
\item $ABC$ مثلث فيه $AB=5$ ، $AC=3$ و $\widehat{A}=\dfrac{2\pi}{3}$ ، إذن $BC$ يساوي:
\ABC{$7$}{$\sqrt{19}$}{$49$}
\item المستقيم الذي يشمل $A(1;1)$ و $\vec n(2;-3)$ ناظمي له معادلته:
\ABC{$3x+2y-5=0$}{$2x-3y+1=0$}{$2x-3y-1=0$}
\end{enumerate}
\end{tamrin}

% ------------------------------------------------- 18
\begin{tamrin}{صحيح أم خطأ}
أجب بصحيح أو خطأ مع التعليل:
\begin{enumerate}
\item إذا كان $\vec u\cdot\vec v=0$ فإن $\vec u=\vec0$ أو $\vec v=\vec0$.
\item $\nV{\vec u+\vec v}^2=\nV{\vec u}^2+\nV{\vec v}^2$ إذا وفقط إذا كان $\vec u\perp\vec v$.
\item المستقيمان $2x-y+1=0$ و $x+2y-3=0$ متعامدان.
\item المعادلة $x^2+y^2-2x+4y+6=0$ هي معادلة دائرة.
\item من أجل كل عدد حقيقي $x$: \ $\cos2x=1-2\cos^2x$.
\item من أجل كل عدد حقيقي $x$: \ $\sqrt3\cos x+\sin x=2\cos\Big(x-\dfrac\pi6\Big)$.
\item إذا كان $\V{AB}\cdot\V{AC}<0$ فإن الزاوية $\widehat{BAC}$ منفرجة.
\item النقطة $A(1;2)$ تنتمي إلى الدائرة التي مركزها $O$ ونصف قطرها $\sqrt5$.
\end{enumerate}
\end{tamrin}

% =====================================================================
\section{تمارين نموذجية}
% ------------------------------------------------- 19
\begin{tamrin}{تمرين نموذجي 1: عناصر مثلث في معلم}
في معلم متعامد ومتجانس $(O;\vec i,\vec j)$ نعتبر النقط: $A(1;4)$ ، $B(-2;-2)$ و $C(4;1)$.
\begin{enumerate}
\item احسب $AB$ ، $BC$ و $AC$ ، ثم استنتج طبيعة المثلث $ABC$.
\item احسب $\V{AB}\cdot\V{AC}$ ، ثم $\cos\widehat{BAC}$ وقيمة مقرّبة لـ $\widehat{BAC}$ بالدرجات.
\item
\begin{sq}
\qa اكتب معادلة للمستقيم $(BC)$.
\qb احسب المسافة بين $A$ والمستقيم $(BC)$ ، ثم استنتج مساحة المثلث $ABC$.
\end{sq}
\item
\begin{sq}
\qa اكتب معادلتي الارتفاعين المتعلقين بالرأسين $A$ و $C$.
\qb استنتج إحداثيي $H$ نقطة تلاقي ارتفاعات المثلث $ABC$.
\end{sq}
\item
\begin{sq}
\qa عيّن إحداثيي $\Omega$ مركز الدائرة $(\mathcal C)$ المحيطة بالمثلث $ABC$ ، ثم اكتب معادلة لها.
\qb ليكن $G$ مركز ثقل المثلث $ABC$. \ بيّن أن $\V{\Omega H}=3\V{\Omega G}$.
\end{sq}
\item اكتب معادلة المماس $(T)$ للدائرة $(\mathcal C)$ في النقطة $A$.
\end{enumerate}
\end{tamrin}

% ------------------------------------------------- 20
\begin{tamrin}{تمرين نموذجي 2: عبارة مثلثية}
من أجل كل عدد حقيقي $x$ نضع: $A(x)=\cos x-\sqrt3\sin x$.
\begin{enumerate}
\item بيّن أنه من أجل كل عدد حقيقي $x$: \ $A(x)=2\cos\Big(x+\dfrac\pi3\Big)$.
\item حل في المجال $[0;2\pi[$ المعادلة $A(x)=1$.
\item عيّن القيمة الحدية العظمى والقيمة الحدية الصغرى لـ $A(x)$ ، وقيم $x$ التي تحققهما.
\item احسب $A\Big(\dfrac\pi4\Big)$ بطريقتين ، ثم استنتج القيمة المضبوطة لـ $\cos\dfrac{7\pi}{12}$.
\end{enumerate}
\end{tamrin}

\end{document}
